\section{Introduction}
\label{sec:introduction}

% =========================================================================
% RL-first introduction, ~1.5 to 2 pages.
% =========================================================================

Approximate $Q$-learning over structured context MDPs offers a principled
way to teach small language models to deliberate
\citep{sutton2018reinforcement, mnih2015human}. The state is a
structured context, the action set is a small collection of cognitive
verbs, the transition model is induced by a frozen LLM, and the reward
is anchored to a verifier that emits a scalar signal on terminal states.
A learned $\Qpsi$ then estimates the long-run expected verifier reward
from a state under each available action, and any of the standard
behavior policies (greedy, softmax, $\epsilon$-greedy) can drive
exploration during sample collection. Compositional reasoning benefits
from being able to decompose problems into independently checkable
sub-problems, because each sub-problem can contribute terminal verifier
signal without depending on the rest. Consider a triangle with vertices
$A$, $B$, $C$ and side lengths $a$, $b$, $c$, and the claim that
$a^2 + b^2 = c^2$. A monolithic verifier asked to confirm this claim
must perform two qualitatively different checks at once: that the
configuration is a right triangle with hypotenuse $c$, and that the
arithmetic relationship between the squared side lengths holds. A
decomposed approach instead splits the parent claim into two
sub-claims, ``the triangle is right-angled at $C$ (by some geometric
fact in the problem)'' and ``the squared side lengths satisfy
$a^2 + b^2 = c^2$ (by arithmetic given the right-angle property),''
each addressing a verifier-checkable sub-problem in isolation. The
geometric premise might be discharged by a perpendicularity check on
two vectors; the arithmetic conclusion might be discharged by a
symbolic algebra system. Neither sub-verifier needs to understand the
other's check, and the approximate $Q$-function (optionally aided by
MCTS as a policy improvement operator) can credit-assign reward up
from the two leaves to the parent claim
\citep{kocsis2006bandit, coulom2006efficient, anthony2017thinking}.
The hierarchical reinforcement learning tradition has long argued that
such structural decomposition is what makes compositional reasoning
tractable in sequential decision-making
\citep{sutton1999between, dietterich2000hierarchical,
vezhnevets2017feudal}.

The risk is that not all recursion is productive. Within the same
schema that permits the Pythagorean decomposition above, an agent may
also produce a critique of a critique of a critique, or a hypothesis
about a hypothesis about a hypothesis, with no verifier-checkable
leaf-level claim at the bottom of the chain. Such trajectories are
structurally legal: nothing in the type system prevents arbitrarily
deep parent-pointer chains, and the verbs that introduce recursive
structure are available whenever a parent entity exists. They are
nonetheless rumination rather than reasoning, in the operational
sense that terminal verifier reward never enters the trajectory at
the leaves and the approximate $Q$-function has no way to assign
positive value through them. Approximate $Q$-learning will eventually
learn to avoid such trajectories, since their Bellman targets remain
at the value-of-non-terminal baseline; but discovering this empirically
costs sample budget, and the time spent ruminating in simulation is
time not spent learning where productive decomposition pays off. The
two modes share an action space, a state schema, and a serialization
template; from syntax alone they are indistinguishable. The question
this paper answers is: when and how should an approximate $Q$-learning
system over an LLM substrate recur, and what structural prior selects
the productive mode without foreclosing the unproductive one's
structural possibility?

\paragraph{Position.}
This is a methodology and position paper, written as a companion to a
framework paper at \texttt{paper/} that establishes the setting. The
framework paper presents verifier-anchored approximate $Q$-learning
over LLM-induced context MDPs for small language model deliberation,
with structured states, a small fixed set of cognitive-verb actions,
a frozen LLM as both encoder and stochastic transition model, MCTS
as an optional sample-efficient policy improvement operator, and a
verifier-anchored reward as the only Bellman target. Its state schema
endows every hypothesis, subgoal, and verifier call with a stable
identifier and with implicit parent references, precisely so that
hypotheses can refine other hypotheses, subgoals can have nested
sub-subgoals, and verifier calls can follow up on earlier ones. The
framework paper notes the forward compatibility, defers exercise of
those fields to a later phase, and gestures briefly at the
productive-versus-unproductive distinction that governs whether
activation will pay off. The present paper takes that gesture
seriously. The design choice deserves its own treatment because its
tradeoffs are non-obvious: the schema that admits productive
decomposition also admits unproductive meta-criticism, and the
architectural job is not to ban one mode and admit the other but
to shape the learning dynamics so that the productive mode dominates
within a finite sample budget. We do not yet report experiments. The
contribution is the conceptual development of the design, together
with a pre-registered ablation grid that makes the central claims
falsifiable.

\paragraph{Contributions.}
We make five contributions:
\begin{enumerate}[leftmargin=*]
  \item A productive-versus-unproductive distinction for structural
        recursion as a property of trajectories under a verifier in
        the framework's approximate $Q$-learning system
        (Section~\ref{sec:method-productive},
        Section~\ref{sec:method-mcts}). Productive recursion is
        structural decomposition with verifier signal at the leaves
        of the implicit tree; unproductive recursion is meta-criticism
        that receives no terminal reward beneath. The distinction is
        operational rather than syntactic: the same action space and
        state schema produce both, the verifier and depth budget
        together determine which trajectories carry learning signal,
        and $\Qpsi$ must learn to prefer the productive regime
        within a finite sample budget.
  \item The depth-limit mechanism as a structural prior on the
        learning system that selects productive recursion without
        foreclosing the unproductive case, framing bounded
        per-hypothesis tree depth $\depthlim$ as a property of the
        environment rather than as a soft reward term
        (Section~\ref{sec:method-depthlimit}). The hard-mask choice
        is justified on framework-paper consistency grounds: a soft
        depth penalty violates the verifier-anchored-reward-only
        principle and admits arbitrarily deep simulation descent
        when MCTS is used as the policy improvement operator. The
        depth limit interacts directly with verifier signal density,
        since a trajectory of depth $d$ produces at most $d$
        verifier-anchored events between root and leaf.
  \item An argument that the framework's seven-verb action set
        $\Aset = \{a_1, \ldots, a_7\}$ is sufficient for arbitrary
        recursive structure via identifier-based targeting, with no
        new categorical actions required
        (Section~\ref{sec:method-targeting}). Targeting actions
        through stable identifiers means
        \textsc{Decompose}\texttt{(hyp\_id)} has the same semantics
        regardless of where \texttt{hyp\_id} sits in the implicit
        tree; the verb-instantiator's free-form identifier choice is
        the only thing that varies. The branching factor at decision
        nodes stays at seven (or fewer, after availability masking),
        and the fixed action-set assumption underlying both $\Qpsi$
        and PUCT-style policy improvement is preserved without
        progressive widening machinery.
  \item A serialization argument that flat tokens with depth markers
        let a frozen LLM resolve recursive tree structure through
        native causal self-attention, with no graph neural network,
        no tree-LSTM, and no recursive cell
        (Section~\ref{sec:method-serialization}). We contrast this
        with the explicit tree-encoder tradition
        \citep{socher2013recursive, tai2015improved,
        velickovic2018gat} and argue that modern LLMs' in-context
        attention is a sufficient substitute, with the
        parent-pointer field and the depth integer providing
        redundant signal for the implicit tree's topology. The
        canonicalization implications of recursive state, and their
        interaction with the policy improvement operator's tree
        reuse, are flagged briefly (Section~\ref{sec:method-canon}).
  \item A pre-registered ablation grid that varies the per-hypothesis
        tree depth limit across
        $\depthlim \in \{0, 3, \infty\}$, holding the framework
        paper's benchmarks, base models, and learning architecture
        fixed (Section~\ref{sec:experimental-plan}). The grid is
        designed to falsify both extremes: $\depthlim = 0$ (no
        recursion, expressivity bottleneck) and
        $\depthlim = \infty$ (unbounded recursion, rumination
        consumes the sample budget). The middle setting
        $\depthlim = 3$ is the predicted productive regime,
        motivated by the natural decomposition depth of
        human-written solutions to competition mathematics
        problems.
\end{enumerate}

\paragraph{Paper structure.}
Section~\ref{sec:background} recaps Markov decision process
foundations, the approximate $Q$-learning view of LLM-induced context
MDPs inherited from the framework paper, and the hierarchical
reinforcement learning and tree-structured-encoder literatures that
the present paper joins and departs from. Section~\ref{sec:method}
is the core: the productive-versus-unproductive distinction with its
formal definition, the depth-limit mechanism with its amended
availability table, action targeting at arbitrary depth, the
serialization treatment, the interaction with MCTS as a policy
improvement operator, and a brief note on the canonicalization
implications. Section~\ref{sec:experimental-plan} specifies the
pre-registered depth-limit ablation grid and the three hypotheses
it commits to. Section~\ref{sec:related-work} positions the
contribution against hierarchical reinforcement learning, recursive
and tree-structured neural networks, and recent verifier-anchored
LLM reasoning systems such as AlphaProof
\citep{deepmind2024alphaproof} and Tree of Thoughts
\citep{yao2023tree}. Section~\ref{sec:discussion} discusses
limitations, threats to validity specific to the recursion claim,
and extensions to adaptive depth limits and DAG-structured
deliberation. Section~\ref{sec:conclusion} concludes.
